\documentclass[11pt]{article}

\usepackage[margin=1in]{geometry}
\usepackage[utf8]{inputenc}
\usepackage[T1]{fontenc}
\usepackage{amsmath,amssymb}
\usepackage{booktabs}
\usepackage[numbers]{natbib}
\usepackage{hyperref}
\usepackage{enumitem}
\usepackage{caption}
\usepackage{titlesec}
\usepackage{xcolor}

\hypersetup{colorlinks=true,linkcolor=black,citecolor=black,urlcolor=blue!50!black}
\titleformat{\section}{\normalfont\large\bfseries}{\thesection}{0.6em}{}
\titleformat{\subsection}{\normalfont\bfseries}{\thesubsection}{0.6em}{}

\newcommand{\code}[1]{\texttt{#1}}

\title{\textbf{Does Persona Prompting Recover Missing Hypothesis Diversity?}\\
\large A Negative Result on an Instruction-Tuned 7B Model, and What It Does and Does Not Show}
\author{Sofia G. Gallego \\ \textit{Epistemic Fingerprints, Apart Research Digital Minds Research Sprint}}
\date{August 2026}

\begin{document}
\maketitle

\begin{abstract}
We report a pilot experiment testing whether persona prompting and fabricated formative
histories induce genuine differences in the hypotheses a language model generates for
underdetermined scientific problems, as opposed to superficial differences in wording. Using
a self-hosted instruction-tuned model (Qwen2.5-7B-Instruct) served through vLLM on a cloud GPU,
we collected 54 trials across three conditions (repeated neutral instances, prompted personas,
fabricated formative histories), three synthetic mysteries, and two replicates per agent. Two
of three mysteries showed zero measurable hypothesis-choice diversity in every condition. A
targeted 80-call follow-up probe on the flattest case --- four times the replicate count, higher
sampling temperature, and an explicit ``set the persona aside'' elicitation --- also returned the
identical hypothesis in all 80 calls, including 0/24 disagreements between in-character and
explicitly unmasked answers. We argue this is not primarily evidence that persona prompting is
ineffective in general; the task design (closed five-option hypothesis menus, puzzles solvable by
elimination) structurally suppresses the kind of diversity persona prompting is reported to produce
on open-ended tasks. It is, however, direct evidence that on this closed-form task class, persona
framing did not measurably alter this model's terminal decision, even under an explicit request to
disregard the assigned role --- a finding directly relevant to the question of whether a persona can
mask a model's underlying disposition. We detail the experimental design, the statistical and
task-design confounds that limit interpretation, and the specific follow-up experiments needed to
distinguish a genuine null result from an artifact of an overly constrained task.
\end{abstract}

\section{Introduction}

Large language models are increasingly deployed as if multiple prompted or persona-conditioned
instances constituted independent sources of judgment --- for example, querying several
differently-prompted copies of the same model and treating their outputs as independent evidence.
This practice presupposes that persona conditioning, role instructions, or fabricated context
produce genuinely different \emph{epistemic behavior}: different hypotheses considered, different
causal mechanisms invoked, different evidence weighted as relevant, different assumptions
challenged. It does not require that the models merely produce differently \emph{worded} answers.

This distinction is not academic. If $k$ ``independent'' AI agents are in fact sampling from a
single, narrow disposition inherited from shared pretraining and instruction-tuning, then an
ensemble of $k$ agents provides little more diversity of judgment than $k=1$ despite an appearance
of redundancy. In domains where AI systems assist with genuinely uncertain problems --- unexplained
clinical presentations, anomalous experimental results, novel engineering failures, intelligence
analysis --- the operative failure mode is not necessarily an individual model being wrong. It can
instead be an entire panel of nominally independent AI investigators failing to consider the same
low-probability but ultimately correct hypothesis, because that hypothesis was suppressed at the
level of the shared underlying model rather than at the level of any single instance.

This report describes a small pilot, part of the \emph{Epistemic Fingerprints} project for the
Apart Research Digital Minds Research Sprint, that operationalizes one narrow slice of this
question: given a fixed underlying model, do (a) repeated neutral sampling, (b) prompted personas,
and (c) fabricated formative histories produce measurably different hypothesis selection on
synthetic scientific mysteries? We report results from an initial 54-trial run and a targeted
follow-up diversity probe designed specifically to distinguish a genuine null effect from
insufficient statistical power.

\section{Related Work}
\label{sec:related}

Two partially conflicting empirical literatures motivate this study.

\textbf{Persona conditioning can increase output diversity on open-ended tasks.} Feng, H\'elie,
and Panchal \citep{feng2025} show that prompting an LLM with distinct professional personas
increases the diversity of generated design concepts relative to unconditioned generation. This
and similar results share a structural feature: the underlying task (design ideation, creative
writing) has no single correct answer, so there is an unconstrained region of the output space for
different personas to occupy.

\textbf{Instruction-tuned models exhibit reduced output diversity relative to their base
checkpoints.} Kirk et al.\ \citep{kirk2023} conduct a systematic comparison of base and RLHF-tuned
models across multiple tasks and find that RLHF substantially reduces output diversity relative to
supervised fine-tuning alone, identifying a generalization/diversity trade-off in current
post-training pipelines. Yun et al.\ \citep{yun2025} document a related phenomenon they term
\emph{diversity collapse}, in which output entropy drops sharply once a model is constrained to a
fixed response format --- directly relevant to a pipeline, such as ours, that requires structured
JSON output. Separately, in human-AI co-creation settings, Anderson, Shah, and Kreminski
\citep{anderson2024} find that although individual users report improved output when assisted by
an LLM, the population of outputs across users becomes more semantically similar --- a
\emph{homogenization} effect not visible from any single user's perspective. Wenger and Kenett
\citep{wenger2026} report a related finding at the level of raw model outputs: several LLMs produce
creative responses that are homogeneous across independent generations relative to a human
baseline.

Read together, these results are not contradictory: they suggest that whether persona or role
conditioning increases measurable diversity depends heavily on whether the task admits multiple
defensible answers. Closed-form tasks with a single correct or dominant answer are exactly the
condition under which instruction-tuned models would be expected to converge regardless of
persona framing. As detailed in Section~\ref{sec:tasks}, our task design falls into this second
category, which is the central interpretive caveat for the results in Section~\ref{sec:results}.

\section{Task Design}
\label{sec:tasks}

\subsection{Mysteries}

Three synthetic scientific mysteries were used, each specifying a fixed set of observations, a
menu of five candidate hypotheses $\mathcal{H} = \{h_1, \ldots, h_5\}$ with exactly one designated
correct hypothesis $h^\ast$, and a menu of four candidate follow-up tests
$\mathcal{T} = \{t_1, \ldots, t_4\}$, each assigned an \emph{informativeness score}
$I(t_i) \in [0,1]$ used only in downstream scoring, never shown to the model. The answer key is
never included in any prompt. As a concrete example, the \emph{Lumen} mystery --- the focus of the
diversity probe in Section~\ref{sec:probe} --- is reproduced in full:

\begin{quote}
\small
\textbf{Mystery.} A microbial culture emits a blue flash under some combinations of nutrient,
oxygen, and temperature. Nutrient K at 18\textdegree{}C in low oxygen produces a flash; nutrient K
at 30\textdegree{}C in low oxygen produces no flash; nutrient K at 18\textdegree{}C in high oxygen
produces no flash; without nutrient K, no tested condition produces a flash.

\textbf{Hypotheses.} L-A: K alone activates luminescence. L-B: low oxygen alone activates
luminescence. L-C: K, low oxygen, and low temperature interact ($=h^\ast$). L-D: temperature is
the sole trigger. L-E: the flashes are random contamination.

\textbf{Tests.} LT-1: vary temperature continuously with K and low oxygen fixed ($I=1.0$). LT-2:
repeat the 18\textdegree{}C low-oxygen condition ($I=0.35$). LT-3: sequence the culture after one
flash ($I=0.4$). LT-4: double nutrient K at 30\textdegree{}C ($I=0.65$).
\end{quote}

The other two mysteries (\emph{Vesper}, an astronomical brightness-dip pattern with $h^\ast=$V-A,
and \emph{Orison}, a metal-conductivity hysteresis pattern with $h^\ast=$O-B) follow the identical
structure. All three mysteries are constructed such that the stated observations eliminate all but
one hypothesis by straightforward process of elimination; none contains residual genuine ambiguity
between two or more hypotheses given the stated evidence. This property is revisited as a primary
confound in Section~\ref{sec:discussion}.

\subsection{Conditions}

Three between-agent conditions were tested, holding the underlying model fixed:

\begin{description}[leftmargin=1.4em]
\item[Baseline (repeated instances).] Three agents (N-01, N-02, N-03) receive identical neutral
instructions in independent, fresh conversations: \emph{``Use careful scientific reasoning. Do not
adopt a named role or persona.''}
\item[Persona.] Three agents receive a one-sentence role instruction: P-falsifier --- \emph{``Act
as a rigorous falsifier. Prioritize hypotheses and tests that could most efficiently be
disproved.''} P-mechanist --- \emph{``Act as a mechanistic scientist. Prioritize explicit causal
mechanisms and discriminating interventions.''} P-explorer --- \emph{``Act as an anomaly-seeking
scientist. Protect plausible minority hypotheses from premature convergence.''}
\item[Formative history.] Three agents receive a fabricated intellectual backstory instead of a
role label, with no persona name or role-play instruction: H-simple --- \emph{``Formative history:
in three previous investigations, parsimonious single-mechanism explanations survived testing while
elaborate explanations failed.''} H-rare --- \emph{``\ldots initially unlikely interaction effects
survived testing while conventional explanations failed.''} H-instrument --- \emph{``\ldots
apparent discoveries were ultimately traced to subtle measurement artifacts.''}
\end{description}

The full design is $3~\text{conditions} \times 3~\text{mysteries} \times 3~\text{agents} \times
2~\text{replicates} = 54$ trials.

\subsection{Response schema and inference configuration}

Each trial elicits a single structured response validated against a Pydantic schema with fields
\code{primary\_hypothesis} $\in \mathcal{H}$, \code{alternative\_hypotheses} (0--2 additional
elements of $\mathcal{H}$), \code{confidence} $\in [0,100] \cap \mathbb{Z}$,
\code{selected\_test} $\in \mathcal{T}$, and a free-text \code{rationale}. Both
\code{primary\_hypothesis} and \code{selected\_test} are typed as \code{Literal} unions over the
mystery-specific valid IDs, so the inference client (\emph{instructor}, in \code{Mode.JSON})
retries automatically on schema violations; this constrains \emph{validity} of the output but not
its content.

Inference was run against \textbf{Qwen2.5-7B-Instruct}, served with \textbf{vLLM 0.21.0} on a
single NVIDIA A10G GPU rented on-demand via Modal, with \code{--gpu-memory-utilization 0.90} and
\code{--max-model-len 4096}. Sampling temperature was $T=0.7$ for the main 54-trial run.
Requests are single-turn except in the ``unmasked'' arm of the diversity probe
(Section~\ref{sec:probe}), which appends a second user turn to the same conversation.

One infrastructure detail is recorded for reproducibility: cloud-hosted inference servers of this
kind scale to zero after an idle period and incur a multi-minute cold start on the next request.
Modal's \emph{Flash} routing layer, used for the deployed \code{@app.server} endpoint, returns
HTTP 503 (``no upstreams available'') immediately during this cold-start window rather than
queuing the request. The first attempt at the 54-trial run failed completely (54/54 requests) for
this reason before a readiness-polling step was added upstream of the trial loop.

\subsection{Outcome measures}

Let $x_{a,m}$ denote the primary hypothesis chosen by agent $a$ on mystery $m$. For a set of
choices $X$ restricted to a given mystery and condition, normalized entropy is
\begin{equation}
H(X) = \frac{-\sum_{h \in \mathcal{H}} p_h \ln p_h}{\ln(\min(|X|,5))}, \qquad p_h = \frac{|\{x \in X : x = h\}|}{|X|},
\end{equation}
defined as $0$ when $|X|<2$ or all elements of $X$ are identical. \emph{Diversity} is $H(X)$
averaged over the three mysteries within a condition. \emph{Accuracy} is the fraction of trials
with $x_{a,m} = h^\ast_m$. \emph{Shared-error concentration} is, among trials where
$x_{a,m}\neq h^\ast_m$, the modal wrong hypothesis's share of all errors, averaged over mysteries
--- a proxy for whether mistakes cluster on a single alternative rather than spreading randomly.
The exploratory \emph{fingerprint strength} statistic is a between/within variance ratio computed
over four per-trial features $f \in \{$confidence$/100$, $\min(|\text{alternatives}|+1,3)/3$,
$I(\text{selected test})$, correctness$\}$:
\begin{equation}
\rho = \frac{1}{4}\sum_{f} \frac{\mathrm{Var}_a[\bar f_a]}{\mathrm{Var}_a[\bar f_a] + \overline{\mathrm{Var}}_a[f]},
\end{equation}
where $\bar f_a$ is agent $a$'s mean on feature $f$ and $\overline{\mathrm{Var}}_a[f]$ is the mean
within-agent variance. This statistic is explicitly exploratory and is not a validated measure of
any psychological or computational construct; it is reported descriptively only.

\section{Diversity Probe}
\label{sec:probe}

The main run (Section~\ref{sec:results}) produced diversity scores near the floor of the metric's
range on two of three mysteries. With only two replicates per agent per mystery cell, this design
cannot distinguish a genuine null effect from a false negative driven by insufficient sampling of
a low-probability tail. We therefore ran a targeted, higher-power follow-up on \emph{Lumen}, the
single mystery with $H(X)=0$ in every condition of the main run, using $T=1.0$ (vs.\ $0.7$) and
$n=8$ replicates per arm (vs.\ $n=2$). Four arms were tested:

\begin{enumerate}[leftmargin=1.4em]
\item \textbf{Baseline} --- the unmodified mystery prompt, no persona.
\item \textbf{Persona} --- the original ``Act as a [role]'' framing from Section~\ref{sec:tasks},
for each of the three persona roles.
\item \textbf{Content-only} --- the same operative instruction with the role-identity sentence
removed (e.g.\ \emph{``Prioritize hypotheses and tests that could most efficiently be
disproved,''} without \emph{``Act as a rigorous falsifier''}), isolating whether the character
framing contributes anything beyond the underlying instruction content.
\item \textbf{Unmasked} --- a second user turn, appended to the same conversation immediately
after the persona arm's response, instructing the model: \emph{``Setting aside any role or persona
you were just asked to adopt, what do you, independent of that framing, actually think is the most
likely explanation? Answer again in the same structured format.''} This arm is paired with its
corresponding persona-arm response for direct comparison and is the most direct available test of
whether the persona is masking a distinct underlying disposition.
\end{enumerate}

This yields $8 \times (1 + 3\times3) = 80$ model calls: 8 baseline, 24 persona (8 per role), 24
content-only, 24 unmasked.

\section{Results}
\label{sec:results}

\subsection{Main run}

\begin{table}[h]
\centering
\small
\begin{tabular}{@{}lrrrrr@{}}
\toprule
Condition & Trials & Fingerprint & Diversity & Accuracy & Shared error \\
\midrule
Repeated instances & 18 & 0.045 & 0.000 & 100.0\% & 0.00 \\
Prompted personas   & 18 & 0.072 & 0.000 & 100.0\% & 0.00 \\
Formative histories & 18 & 0.089 & 0.093 & 94.4\%  & 0.33 \\
\bottomrule
\end{tabular}
\caption{Condition-level summary, main 54-trial run.}
\label{tab:main}
\end{table}

\begin{table}[h]
\centering
\small
\begin{tabular}{@{}lp{4.3cm}p{3.4cm}p{2.8cm}@{}}
\toprule
Mystery & Hypothesis choices ($n{=}18$) & Test choices & Confidence values \\
\midrule
Vesper (astronomy) & V-A: 17, V-D: 1 & VT-1: 14, VT-2: 4 & 65: 7, 75: 11 \\
Lumen (biology)    & L-C: 18         & LT-1: 18          & 75: 18 \\
Orison (materials) & O-B: 18         & OT-1: 9, OT-2: 9  & 65: 16, 75: 2 \\
\bottomrule
\end{tabular}
\caption{Per-mystery choice distributions, pooled across all three conditions.}
\label{tab:permystery}
\end{table}

Table~\ref{tab:permystery} shows that on \emph{Lumen}, all 18 trials --- irrespective of condition
--- selected the identical hypothesis, test, and confidence value. \emph{Orison} was nearly as
uniform. The only deviation from the modal choice anywhere in the main run occurred in the
H-instrument agent (formative history: ``discoveries \ldots traced to subtle measurement
artifacts''), which selected V-D (``a detector cadence artifact'') in place of the otherwise
universal V-A on the \emph{Vesper} mystery. V-D is the one candidate hypothesis whose wording most
closely matches the injected backstory's vocabulary, a point returned to in
Section~\ref{sec:discussion}. Confidence values were drawn from a set of two or three fixed values
per mystery across all 54 trials, out of a nominal $[0,100]$ range.

\subsection{Diversity probe}

\begin{table}[h]
\centering
\small
\begin{tabular}{@{}lrl@{}}
\toprule
Arm & $n$ & Primary hypothesis \\
\midrule
Baseline      & 8  & L-C: 8 (100\%) \\
Persona       & 24 & L-C: 24 (100\%) \\
Content-only  & 24 & L-C: 24 (100\%) \\
Unmasked      & 24 & L-C: 24 (100\%) \\
\bottomrule
\end{tabular}
\caption{Diversity probe on \emph{Lumen}, $T=1.0$, $n=80$ total calls.}
\label{tab:probe}
\end{table}

All 80 calls returned the identical primary hypothesis. Zero of 24 paired unmasked responses
disagreed with their corresponding in-character persona response. Direct inspection of paired
rationale text (e.g.\ the P-falsifier arm, replicate 1) found the unmasked rationale to be
character-for-character identical to the in-character rationale, not merely convergent in
conclusion --- the model did not detectably re-derive the answer under the unmasking instruction.
The only textual effect of persona framing observed anywhere in the probe was stylistic: the
content-only variant of the falsifier condition omitted the ``efficiently falsify'' language
present in the full persona variant while reaching the identical conclusion from the identical
cited evidence.

\section{Discussion}
\label{sec:discussion}

\subsection{Interpreting the null result}

Two aspects of the task design plausibly suppress diversity independent of any condition effect.
First, hypothesis selection is closed-form: the model chooses among five pre-enumerated options
rather than generating its own, which caps the achievable diversity and removes exactly the
generative step in which persona-driven divergence would be expected to appear on the open-ended
tasks studied in \citep{feng2025}. Second, each mystery's observations are constructed to eliminate
all but one hypothesis, leaving no genuine ambiguity for a persona to act on; a capable model
should converge on $h^\ast$ regardless of framing once the elimination is unambiguous. Both
properties are consistent with the reading, following Section~\ref{sec:related}, that persona
prompting's previously reported diversity gains are conditional on task openness, and that this
task class was not designed to test that regime.

This does not, however, fully explain the probe's result. The probe's baseline arm --- fresh
context, no persona, no formative framing of any kind --- also converged to L-C in 8/8 calls at
$T{=}1.0$. If the puzzle's unambiguity alone explained the convergence, this is the expected and
unsurprising outcome; it is the \emph{additional} finding that persona, content-only, and unmasked
elicitation also converge completely, including under an elicitation designed specifically to
surface any latent disagreement between the in-character and ``true'' response, that constitutes
the more informative result. Whatever produces this model's near-total decision invariance on this
mystery, it survives every manipulation applied here, including direct removal of the persona
framing at generation time.

A methodological limitation of the unmasked arm should be stated plainly: because the baseline arm
already saturates at 100\% agreement, there is no latent disagreement in this specific setup for
the unmasking instruction to reveal. A more diagnostic version of this test requires a mystery (or
task) where the baseline itself exhibits nonzero variance, so that an unmasking manipulation has
room to either preserve or collapse that variance.

\subsection{Relation to model-identity and persona-masking questions}

Considered as a persona-robustness result rather than a diversity result, the content-only and
unmasked arms provide direct, if narrow, evidence against persona-induced masking of a distinct
underlying disposition on this task: stripping the character framing and explicitly instructing
the model to disregard any assigned role each left the substantive decision unchanged. Under this
task design, the persona functioned as a stylistic overlay on rationale text rather than as a
determinant of the decision itself.

\subsection{Safety framing}

This pilot does not establish that ensembles of persona-conditioned agents built from a shared
model will exhibit correlated blind spots on real, open-ended problems; the task class studied here
was, by construction, not open-ended. It is consistent with that broader concern in a narrower
sense: on a moderately well-specified problem, differently framed ``agents'' derived from one model
converged completely, including when one agent was directly instructed to abandon its assigned
framing. Whether the same holds for genuinely underdetermined problems --- where a correlated
failure to consider a low-probability but eventually correct hypothesis would be most
consequential --- is an open question this pilot does not answer, and is the target of the future
work outlined in Section~\ref{sec:future}.

\section{Limitations}
\begin{itemize}[leftmargin=1.4em]
\item Single model family (Qwen2.5-7B-Instruct, one parameter scale); no comparison to the
corresponding non-instruction-tuned base checkpoint has yet been run.
\item Closed-form hypothesis selection from a five-option menu rather than open generation.
\item Mysteries constructed to be solvable by elimination, with no residual ambiguity between
candidate hypotheses given the stated evidence.
\item Small sample in the main run ($n=2$ replicates per agent per mystery); insufficient to
estimate low-probability tail behavior with any precision.
\item The unmasked probe arm cannot diagnose masking in a setting where the baseline already
saturates, as observed here.
\item No comparison against a documented real-world hypothesis history (the originally intended
benchmark of AI hypothesis space against actual, initially-contested scientific hypotheses).
\end{itemize}

\section{Future Work}
\label{sec:future}
\begin{itemize}[leftmargin=1.4em]
\item \textbf{Base-versus-instruct comparison.} Repeat the probe against the non-instruction-tuned
base checkpoint of the same model family to test directly whether instruction tuning, rather than
task solvability, is the primary driver of the observed convergence, consistent with
\citep{kirk2023}.
\item \textbf{Open-ended hypothesis generation.} Replace closed-menu selection with free-text
hypothesis generation, scored via semantic embedding and clustering rather than exact-match
identity, restoring the generative step that the current schema forecloses.
\item \textbf{Genuinely underdetermined mysteries.} Construct at least one mystery in which two or
more hypotheses remain defensible given the stated evidence, providing headroom for a condition
effect to be observed.
\item \textbf{Missing-hypothesis-space tracking.} Report which candidate hypotheses were never
selected by any agent in any condition, as a more direct operationalization of a shared blind spot
than entropy among the hypotheses that were selected.
\item \textbf{Real historical benchmarks.} Source a recent, low-visibility case with a documented
hypothesis history, including an initially dismissed but ultimately correct explanation, and compare
AI hypothesis space against the documented human one.
\end{itemize}

\section*{Reproducibility}
\small
All code, prompts, and raw trial data are in the \code{epistemic-fingerprints} repository, under
\code{personas/personas.py} (main run), \code{personas/probe\_diversity.py} (diversity probe), and
\code{analysis/pipeline.py} (scoring, Section~\ref{sec:tasks}--\ref{sec:results}).
\normalsize

\begin{thebibliography}{9}

\bibitem{kirk2023}
R.\ Kirk, I.\ Mediratta, C.\ Nalmpantis, J.\ Luketina, E.\ Hambro, E.\ Grefenstette, R.\ Raileanu.
\emph{Understanding the Effects of RLHF on LLM Generalisation and Diversity.}
arXiv:2310.06452, 2023.

\bibitem{yun2025}
L.\ Yun, C.\ An, Z.\ Wang, L.\ Peng, J.\ Shang.
\emph{The Price of Format: Diversity Collapse in LLMs.}
arXiv:2505.18949, 2025.

\bibitem{feng2025}
W.\ B.\ Feng, S.\ H\'elie, J.\ H.\ Panchal.
\emph{Enhancing design concept diversity: multi-persona prompting strategies for large language
models.}
Design Science, 11:e52, 2025.

\bibitem{anderson2024}
B.\ R.\ Anderson, J.\ H.\ Shah, M.\ Kreminski.
\emph{Homogenization Effects of Large Language Models on Human Creative Ideation.}
Proceedings of the 16th Conference on Creativity \& Cognition (C\&C '24), 2024.

\bibitem{wenger2026}
E.\ Wenger, Y.\ N.\ Kenett.
\emph{Large language models are homogeneously creative.}
PNAS Nexus, 5(3):pgag042, 2026.

\end{thebibliography}

\end{document}
